\section{Persian Traditions Chap I – General, part a}

Many people have noted the parallels and similarities between Persia and Spain. Wilfrid Blunt (1) says that the Persians are the Spaniards of the Middle East and that both Spaniards and Persians are proud, hospitable and procrastinating. Sr. Eugenio Montes, in an article in the Madrid newspaper ABC in honor of the 2500th anniversary of the founding of the Persian Empire, says: ``Persia is the Spain of Asia, Spain is the Persia of the West.'' Reading this, one friend commented that it is no coincidence that Eugenio Montes is a Gallego, i.e., from Galicia. We will have more to say of this later.

What, one may ask, would cause such distinguished gentlemen as Mr. Blunt and Sr. Montes to affirm such close relations between two countries so distant one from the other? There are many reasons, some of which neither of these two distinguished men of letters seem to be aware.

The geographical similarities between the Iberian Penninsula and the Iranian Plateau are obvious, though of course the Iranian Plateau is on a much larger scale. Both are very largely cut off from the continents of which they form a part by high, rugged mountain ranges; both have a high, cold and somewhat barren central plateau, broken in places by mountains. In the province of Mazanderan one has the Iranian equivalent of the lush, green Cantabrian Coast (even the typical barns of Mazanderan resemble those of Asturias and Santander) in Persian Azerbaijan the equivalent of the wheat plains of Leon. In the Persian Gulf Coast, one has the equivalent of the Southern Mediterranean Coast of Spain. The resemblances are not purely geographical; there are also profound ethnic and historical parallels and similarities.

Spain has a very strong Celtic heritage. At one time or another the Celts penetrated the entire Penninsula and strongly occupied its Western 75 per cent. The Celts were of course never expelled and in spite of Roman, Germanic and Semitic superstrata left a rich heritage to their descendants. In Spain Celtic influences are visible in physical types (particularly in Asturias, Leon and Old Castile), music, artistic and literary forms, festivals, folklore, spiritual values, temperament and general character traits. As we shall see, in Medieval Spain, both Christian and Muslim, poets composed verse using versification forms of Celtic origin, though writing in Latin, Romance (by ``Romance'' I refer to the language called in Arabic \emph{Lisan} \emph{al}--\emph{Ajjam}, i.e., ``the non-Arabic Language''), Classical Arabic, Vulgar Arabic, Hebrew, Catalan, Gallego-Portuguese, Aragonese and Castilian, and sang them to melodies based on Celtic musical modes. No one who has any knowledge of Celtic Studies has the slightest doubt as to the truth of what I am saying, not anyone who has any first-hand knowledge of Spain on the one hand and Celtic countries such as Ireland, Scotland, Wales and Brittany on the other. Celtic elements are present in the phonetics, grammar and vocabulary of all the Romance languages of the Penninsula. So powerful is the Celtic element in Spain that many Spaniards of my acquaintance feel that Spain should be considered as one of the Celtic countries along with Ireland, Scotland, Wales and Brittany, arguing that a Spaniard does not cease to be a Celt because he speaks a Romance language any more than an Irishman ceases to be a Celt because he speaks English. It should never be forgotten that Spain is a land of heather as well as olive trees, a land of bagpipes as well as guitars, and that this is true of Muslim Spain as well as Christian Spain.

Before we go any further, some things need to be clarified. In this work I will use the term ``Aryan'' with a certain frequency, so let us define our terms, Firstly, remember your Aristotle; all Aryans are Indo-Europeans, but not all Indo-Europeans are Aryans. The term ``Aryan'' in reality can be accurately applied to three, and only three Indo-European peoples, i.e., the Indo-Aryans, the Iranians, and the Celts. The word \emph{Arya} or \emph{Aryan} is Sanskrit. The name \emph{Iran} means ``Land of the Aryans'', as does \emph{Erinn}, the native Celtic word for ``Ireland'' (the name \emph{Ireland} is a Viking word), and name ``Aryan'' is found in numerous ancient Celtic place names and tribal names, as we shall demonstrate. Finally, as we shall also demonstrate, the Indo-Aryans, Iranians and Celts have so many special affinities in common to warrant lumping them together as a sub-grouping within the larger Indo-European context. When I use the word or name ``Aryan'', I refer \emph{only} to the Indo-Aryans, Iranians and Celts, \emph{never} to the Germanic peoples. The ancient Germanic peoples did not know the name or word ``Aryan'', and never called themselves by this name.

Hitler was both an ignoramus and a congenital liar, who claimed that the Germanic peoples are ``pure Aryans'', when, in reality, they are not Aryans at all; the Germanic peoples are no more ``pure Aryans'' than they are ``pure Albanians'' (Albanians are Indo-Europeans, but not Aryans). Many people who proclaim themselves to be ``intellectuals'' use the word ``Aryan'' as a synonym for the Germanic peoples, as did Hitler; in this they are showing themselves to be ``pseudo-intellectuals''. Immediately some of these pseudo-intellectuals will try to defend themselves by saying that they are not knowledgeable in Indo-European studies, that it is not their field. Now, there is a word for people who try to pontificate on that of which they know nothing: that word is ``fool''. There is a Cajun expression for people of this sort: ``alligators''. Why? Because an alligator has an enormous mouth and a very tiny brain. Finally, as we shall see, the living language closest to the original Indo-European language is Lithuanian, while the country with the largest percentage of natural blondes is Lithuania; but the Lithuanians are not Germanic. So, from every possible point of view, Hitler's claim that the Germanic peoples are ``pure Aryans'', or even ``Aryans'' at all, is arrant nonsense, the ravings of a liar and an ignorant fool.

Those self-proclaimed ``intellectuals'' who continue to use the word ``Aryan'' in the sense in which it was used by Hitler not only show themselves to be ignorant louts, but also give Hitler credit for an intellectual acumen which he totally lacked. Perhaps this is not surprising.

Hitler was very much a ``progressive'', a man of the extreme left, a true son of the so-called ``Enlightenment''. One could hardly find a more ``leftist'' or ``progressive'' name than ``National Socialist German Workers' Party''. Yes, Hitler was very much a socialist; he repeatedly said ``I am a fanatical socialist'', and praised the advantages of a command economy in contrast to private enterprise. Hitler was also anti-Catholic, anti-aristocratic, anti-monarchist, and was very big indeed on ``separation of church and state''. By at least indirectly attributing to Hitler an intellectual acumen of which he was totally lacking, these ``progressive'' pseudo-intellectuals also reveal their secret love and admiration for Hitler.

In summary, whenever I use the word ``Aryan'', I refer \emph{only} to the Indo-Aryans, the Iranians and the Celts, \emph{never} to the Germanic peoples.

There is a song about Ireland, commercial, not traditional, whose last strophe says:

\begin{verse}
So they (the angels) sprinkled it with stardust\\
Just to make the shamrocks grow\\
`Tis the only place you'll find them\\
No matter where you go\\
Then they sprinkled it with silver\\
Just to make the lakes so grand\\
And when they had it finished\\
Sure they called it Ireland.
\end{verse}

Now, as we said above, the native Celtic name of Ireland is \emph{Erinn}, more modern Gaelic \textbf{\emph{Erin}}, while the name ``Ireland'' is a Viking word. Now, to a true Aryan, Celtic son of the ancient Erinn, the idea that the angels would give a Viking name to the Holy Emerald Isle is simply intolerable, as is anything which could conceivably seem to suggest that the Vikings were angels. So, the above song should be boycotted by all sons of the ancient Erinn, as well as all Celts and those sympathetic to the Celts, who consider themselves to be at least partly Celtic by blood and heritage, which includes the majority of the Spanish and Portuguese, as well as a great many Frenchmen. The angels are not Vikings, nor would they give a Viking name to the ancient Erinn, the Land of the Aryans. And the Vikings, being Germanic, were NOT Aryans.

One of the great enigmas of Indo-European studies is the precise nature of the relation between the Celts and the Indo-Iranian peoples. Specialists in the fields of religion, art, customs and literary forms tend to consider the Celts as an Eastern Aryan people migrated to the West, while those whose specialty is more strictly linguistic tend to affirm that the Celts are a Western people akin to the Italic peoples. There are excellent reasons to question this.

For many years it has been generally believed in scholarly circles that the Celtic languages are most closely related to the Italic languages. This has now been put into grave doubt, in fact it now appears to be almost certainly in error. We shall deal with this at some length, as it is of crucial importance to our study.

Crucial to our thesis are the Illyrians, who, together with the Thracians and Dacians, form the ``bedrock'' of the population of the northern part of the Balkan Peninnsula and present-day Romania.

Says Hans Krahe:

\begin{quotex}
The Illyrian languages, which at one time were spoken in a vast area from the Baltic Sea to the Mediterranean and from Western Europe to Asia Minor, something which has been proven by a study of place names. The Illyrian languages achieved such a vast extension thanks to a great movenemnt which began in the 13th century BC, and which is commonly referred to as the ``Aegean migration'', and whose principal impulse was that of the Illyrians themselves. The linguistic remains of the Illyrian language which have come down to us are meagre, compared with that which one might expect from its ancient diffusion. Of the two Illyrian dialects spoken in the Appenine Penninsula (Italy), i.e., Mesapian in Apulia and Calabria and Veneto in northeastern Italy (from whence comes the name Venezia, i.e., ``Venice''), several hundred inscriptions of each have come down to us. To these must be added one brief inscription, the only examp0le of the Illyrian language spoken in the (northwestern) Balkan Penninsula. In addition, there survive a great many place names and personal names, and numerous glosses. (2)
\end{quotex}

Some affirm that the present-day Romanian language, though a Romance language with a Slavic admixture, contains many elements from the Illyrian language and also Dacian or Thracian. Some also believe that the Bulgarian language, though most certainly Slavic -- indeed, Old Bulgarian is the basis of Church Slavonic, the most prestigious of all Slavic languages, see Chapter 8 -- contains Thracian elements. After all, the Illyrians and Dacians or Thracians are the genetic bedrock of the peoples of the northern part of the Balkan Penninsula and present day Romania.


\flrightit{Posted on 2020-03-28 by Michael McClain}

